\begin{lemma}[Right invariance of a shared internal regular bond]%
    \label{lem:peps_regular_region_bond_right_invariance}
    \lean{TNLean.PEPS.regularRegionHalfEdgeRightMul,
        TNLean.PEPS.regularProjectorTwistedRegionMatrix_halfEdgeRightMul,
        TNLean.PEPS.regularProjectorOpenRegionMatrix_halfEdgeRightMul,
        TNLean.PEPS.regularProjectorOpenRegionMatrix_internalEdgeRightMul,
        TNLean.PEPS.regionPhysicalMap_halfEdgeRightMul_of_mem_regularGroundSpace,
        TNLean.PEPS.regionPhysicalMap_slice_halfEdgeRightMul_of_parent_annihilates}
    \leanok
    \uses{def:peps_region_ground_space, def:peps_region_parent_interaction}
    Consider the actual canonical regular contraction on any
    finite graph region $R$. Choose group elements $r_e$ on its
    bonds, with $r_e=1$ on every crossing bond. Simultaneously
    multiply both exposed half-edge labels of each internal bond
    $e$ on the right by $r_e$. Every canonical regional column
    remains unchanged. In particular this holds for one internal
    bond, with both its endpoints in $R$. The same statement
    holds in the presence of arbitrary inserted left group multiplications.

    Let the physical coefficients be $a_v$, and suppose that local
    inverse coefficients $F_v$ satisfy
    \[
        \sum_s F_v(\alpha,s)a_v(\eta,s)
        =\Pi_v(\alpha,\eta),
    \]
    where $\Pi_v$ is the regular averaging projector. Then every
    original regional ground vector exposed by the product of
    the $F_v$ has this internal right invariance. Consequently,
    every complementary slice of an arbitrary global vector
    satisfying the original regional parent condition has this
    invariance after exposure. Such inverse coefficients exist
    for regular $G$-injective tensors.

    No $G$-isometry, flat-connection reconstruction, or global
    closure-span representation is assumed. This is a local
    identity of the contractions appearing in the accessible
    coordinate argument of \cite{Schuch2010PEPS}.
\end{lemma}

\begin{proof}\leanok
    The regular averaging kernel is unchanged when its two
    configurations are multiplied on the right, leg by leg, by
    the same labels: common left translation commutes with these
    right multiplications. Change each summed shared bond label
    $\eta_e$ to $\eta_e r_e$. This is a bijection of the summation
    variables. It fixes the crossing condition, and associativity
    gives the same identity in the presence of inserted left
    group multiplications. The canonical column equality follows.

    Apply the stated inverse coefficients to an actual original
    regional column. The result is its canonical column, so it
    inherits the identity. Linearity extends it to the entire
    original regional ground space. A regional parent constraint
    places every complementary physical slice in that space.
\end{proof}
